\chapter{Neural Networks}

Neural networks are computational models inspired by the human brain, designed to recognize patterns and solve complex problems. They consist of interconnected layers of artificial neurons that process and transmit information.

\section{Artificial Neuron}

An artificial neuron is the basic building block of a neural network. It takes multiple inputs and, after processing them, produces an output. The processing typically involves a weighted sum of the inputs followed by an activation function.
\begin{enumerate}
    \item \ul{Inputs}: The neuron receives multiple inputs, each associated with a weight that determines its importance.
    \item \ul{Combination of Inputs}: The neuron computes a weighted sum of the inputs, often adding a bias term to allow for more flexibility in the output. This is usually called ``net input'':
    \begin{equation*}
        \sum_{i=1}^{n} w_i x_i + b
    \end{equation*}
    \item \ul{Activation Function}: The result of the weighted sum is passed through an activation function, which introduces non-linearity into the model. This allows the network to learn complex patterns.
    \item \ul{Output}: The output of the activation function is the final output of the neuron.
    \begin{equation*}
        y = \varphi\left(\sum_{i=1}^{n} w_i x_i + b\right)
    \end{equation*}
\end{enumerate}

\subsection{Activation Function}

Activation functions are crucial in neural networks as they allow the model to capture non-linear relationships. They are called ``activation'' functions because they determine whether a neuron should be activated or not based on the input it receives. Common activation functions include:
\begin{itemize}
    \item \ul{Identity}: No real activation function. It is really used in the output layer for regression problems where the output can take any real value.
    \begin{equation*}
        \varphi(x) = x
    \end{equation*}
    \item \ul{Step Function}: A binary function that outputs 1 or 0 depending on whether the input is above or below a certain threshold (often 0). Only the positive inputs activate the neuron.
    \begin{equation*}
        \varphi(x) = \begin{cases}
            1 & \text{if } x \geq 0 \\
            0 & \text{if } x < 0
        \end{cases}
    \end{equation*}
    \item \ul{Sign Function}: Similar to the step function but outputs $-1$ or $1$ depending on the sign of the input.
    \begin{equation*}
        \varphi(x) = \sgn(x) = \begin{cases}
            1 & \text{if } x > 0 \\
            -1 & \text{if } x < 0
        \end{cases}
    \end{equation*}

    \item \ul{Logistic/Sigmoid Function}: A smooth, S-shaped function that maps any real-valued number into the $\left]0, 1\right[$ interval. It is commonly used in binary classification problems.
    \begin{equation*}
        \varphi(x) = \frac{1}{1 + e^{-x}}
    \end{equation*}
    \item \ul{Hyperbolic Tangent Function}: Similar to the logistic function but maps inputs to the $\left]-1, 1\right[$ interval. It is often used in hidden layers of neural networks.
    \begin{equation*}
        \varphi(x) = \tanh(x) = \frac{e^{2x} - 1}{e^{2x} + 1}
    \end{equation*}
    \item \ul{Softmax Function}: Used in multi-class classification problems. A parameter $T$ (temperature) can be introduced to control the smoothness of the output distribution.
    \begin{equation*}
        \varphi(x_i) = \frac{e^{\nicefrac{x_i}{T}}}{\sum_{j} e^{\nicefrac{x_j}{T}}}
    \end{equation*}
    where $j$ ranges over all classes. It is important to note that the softmax function is typically applied to the output layer of a neural network when performing multi-class classification, as it converts the raw output scores (logits) into probabilities that sum to 1 across all classes.

    \item \ul{Rectified Linear Unit (ReLU)}: A popular activation function in deep learning.
    \begin{equation*}
        \varphi(x) = \max(0, x) = \begin{cases}
            0 & \text{if } x < 0 \\
            x & \text{if } x \geq 0
        \end{cases}
    \end{equation*}

    It is really useful in hidden layers of deep networks because it helps to mitigate the vanishing gradient problem, allowing for faster training and better performance. Some modern variants of ReLU include:
    \begin{itemize}
        \item \ul{Parametric ReLU (PReLU)}: Introduces a parameter $\alpha$ to allow for a small slope in the negative part of the function.
        \begin{equation*}
            \varphi(x) = \begin{cases}
                \alpha x & \text{if } x < 0 \\
                x & \text{if } x \geq 0
            \end{cases}
        \end{equation*}
        \item \ul{Softplus}: A smooth approximation of ReLU.
        \begin{equation*}
            \varphi(x) = \ln(1 + e^x)
        \end{equation*}
    \end{itemize}

    \item \ul{Exponential Linear Unit (ELU)}: Similar to ReLU but allows for negative outputs, which can help with learning.
    \begin{equation*}
        \varphi(x) = \begin{cases}
            \alpha (e^x - 1) & \text{if } x < 0 \\
            x & \text{if } x \geq 0
        \end{cases}
    \end{equation*}
\end{itemize}

Appart from these, there are many other activation functions that have been proposed in the literature, each with its own advantages and disadvantages depending on the specific problem being addressed.

\section{Neural Networks}

A neural network is a collection of interconnected artificial neurons organized in layers. There are typically three types of layers in a neural network:
\begin{itemize}
    \item \ul{Input Layer}: The layer that receives the initial data.
    \item \ul{Hidden Layers}: One or more layers that process the inputs from the previous layer and pass the results to the next layer. These layers are called ``hidden'' because they are not directly exposed to the input or output.
    
    Deep networks are simply neural networks with many hidden layers, which allow them to learn more complex representations of the data.
    \item \ul{Output Layer}: The layer that produces the final output of the network.
\end{itemize}

In order for a neural network to learn from data and work effectively, it needs to be trained. The following steps are typically involved in the training process:
\begin{enumerate}
    \item \ul{Architecture Design}: The structure of the network is defined, including the number of layers, the number of neurons in each layer, the choice of activation functions and, specially, the network topology (how the layers are connected).
    \item \ul{Initialization}: The weights and biases of the network are initialized, often with small random values.
    \item \ul{Propagation}: The input data is fed through the network, and the output is computed.
    \item \ul{Loss Calculation}: A loss function is used to measure the difference between the predicted output and the actual target values.
    \item \ul{Training}: The error is propagated back through the network to update the weights and biases in a way that minimizes the loss.
    \item \ul{Iteration}: The process is repeated for multiple epochs (steps) until the network converges to a solution or reaches a predefined number of iterations.
    \item \ul{Evaluation}: The performance of the trained network is evaluated using a separate test dataset to ensure that it generalizes well to unseen data.
\end{enumerate}

Given the importance of the training phase, we will focus deeper on it.
\subsection{Training Process of Neural Networks}

The training process of neural networks is crucial for their performance. It has two main phases: the error propagation phase and the weight update phase.
\subsubsection{Backpropagation Phase}
In this phase, the error is propagated backward through the network to compute the error made by each neuron. The error for the output layer is calculated separately (usually using a loss function as the MAE) and then propagated back to the hidden layers. The error for a neuron is computed as follows:
\begin{equation*}
    e_j = \sum_{k} \dfrac{w_{jk}}{\sum\limits_{i} w_{ik}} \cdot e_k
\end{equation*}
where $k$ ranges over the neurons in the next layer, $w_{jk}$ is the weight connecting the current neuron to the $k$-th neuron in the next layer, $i$ ranges over the neurons in the current layer, and $e_k$ is the error of the $k$-th neuron in the next layer. This process continues until the error has been propagated back to all layers of the network.

However, an important problem with backpropagation is the vanishing gradient problem, which occurs when the gradients become very small as they are propagated back through the layers, making it difficult for the network to learn effectively. This is particularly problematic in deep networks with many layers. To mitigate this issue, various techniques have been developed, such as using different activation functions (like ReLU) or implementing skip connections.

\subsubsection{Weight Update Phase}
In this phase, the weights and biases of the network are updated based on the computed errors. It is done as follows:
\begin{equation*}
    w_{ij,\text{new}} = w_{ij,\text{old}} - \eta \cdot \Delta w_{ij}
\end{equation*}
where $\eta$ is the learning rate, and $\Delta w_{ij}$ is the change in the weight $w_{ij}$, which is the weight connecting the $i$-th neuron in the previous layer to the $j$-th neuron in the current layer. The specific method for calculating $\Delta w_{ij}$ can vary, and there are several algorithms for updating the weights, including:
\begin{itemize}
    \item \ul{Gradient Descent}: A general optimization algorithm that minimizes the loss function by iteratively moving in the direction of the steepest descent as defined by the negative of the gradient.
     \begin{equation*}
        \Delta w_{ij} = \dfrac{\partial L}{\partial w_{ij}}
    \end{equation*}
        where $L$ is the loss function. Using the chain rule, this can in practice be computed as:
    \begin{equation*}
        \Delta w_{ij} = o_i\cdot \delta_j
    \end{equation*}
    where:
    \begin{itemize}
        \item $o_i$ is the output of the $i$-th neuron in the previous layer.
        \item $\delta_j$ is the error term for the $j$-th neuron in the current layer, which can be computed as:
        \begin{equation*}
            \delta_j = e_j \cdot \varphi'(\text{net}_j)
        \end{equation*}
        where $e_j$ is the error for the $j$-th neuron and $\varphi'(\text{net}_j)$ is the derivative of the activation function evaluated at the net input of the $j$-th neuron.
    \end{itemize}
    \item \ul{Hebb's Rule}: A simple learning rule based on the principle that neurons that fire together, wire together. The change in weight is proportional to the product of the input and output of the neuron.
    \begin{equation*}
        \Delta w_{ij} = o_i \cdot o_j
    \end{equation*}
    where $o_i$ is the output of the $i$-th neuron in the previous layer and $o_j$ is the output of the $j$-th neuron in the current layer.
    \ul{Stochastic Gradient Descent (SGD)}: A variant of gradient descent that updates the weights using a single training example (or a small batch of examples) at a time, rather than the entire dataset. This can lead to faster convergence and better generalization, especially in large datasets, as it allows the model to escape local minima and explore the loss landscape more effectively.
\end{itemize}

There are however many other algorithms for updating the weights, such as Momentum, RMSProp, Adam, etc., each with its own advantages and disadvantages depending on the specific problem being addressed.
\begin{observacion}
    Optimizers are usually used, which are algorithms that adjust the learning rate and other hyperparameters during training to improve convergence and performance. Some popular optimizers include Adam, RMSProp, and Adagrad.
\end{observacion}


\section{Recurrent Neural Networks (RNNs)}

Recurrent Neural Networks (RNNs) are a type of neural network designed to handle sequential data. They are particularly effective for tasks such as natural language processing, time series analysis, and speech recognition. They learn these temporal dependencies by maintaining a hidden state that is updated at each time step, allowing them to capture information from previous inputs in the sequence. The architecture of RNNs includes loops that allow information to persist, making them suitable for tasks where the order of the data is important. Three types of loops are commonly used in RNNs:
\begin{itemize}
    \item \ul{Direct Connections}: Connecting neurons to themselves, allowing them to retain information from previous time steps.
    \item \ul{Lateral Connections}: Connecting neurons within the same layer, enabling them to share information across time steps.
    \item \ul{Indirect Connections}: Connecting neurons to neurons in previous layers, allowing for more complex interactions between different time steps.
\end{itemize}


% There are two main types of RNNs:
% \begin{itemize}
%     \item \ul{Elman Networks}: A type of RNN where the hidden layer is connected to itself, allowing it to maintain a state that can capture information from previous time steps.
%     \item \ul{Jordan Networks}: A type of RNN where the output layer is connected back to the hidden layer, allowing the network to use its own output as input for the next time step.
% \end{itemize}

\subsection{One-Directional RNNs}
In one-directional RNNs, the information flows in one direction, from the input layer to the output layer. This means that the network can only capture dependencies from previous time steps, making it less effective for tasks that require understanding of future context. One-directional RNNs are commonly used for tasks such as language modeling and time series prediction, where the goal is to predict the next word in a sentence or the next value in a sequence based on the previous context.

\subsection{Bidirectional RNNs}
Bidirectional RNNs are designed to capture dependencies in both directions, allowing them to understand the context from both past and future time steps. This is achieved by having two separate hidden layers: one that processes the input sequence in the forward direction and another that processes it in the backward direction. The outputs from both layers are then combined to produce the final output.

\subsection{Long Short-Term Memory (LSTM) Networks}

Given the vanishing gradient problem, RNNs can struggle to learn long-term dependencies in sequences. To address this issue, Long Short-Term Memory (LSTM) networks were introduced. LSTMs are a type of RNN that includes special units called memory cells, which can maintain information for long periods of time. These cells are controlled by four gates: 
\begin{itemize}
    \item \ul{Input Gate}: Controls how much of the new information should be stored in the memory cell.
    \item \ul{Forget Gate}: Controls how much of the existing information in the memory cell should be retained or discarded.
    \item \ul{Output Gate}: Controls how much of the information in the memory cell should be output to the next layer.
    \item \ul{Candidate Gate}: Generates new candidate values that could be added to the memory cell.
\end{itemize}

The LSTM architecture allows it to learn long-term dependencies in sequences, making it effective for tasks such as language modeling, machine translation, and speech recognition. By controlling the flow of information through the gates, LSTMs can selectively retain or discard information, allowing them to capture complex patterns in sequential data.

\subsubsection{Gated Recurrent Units (GRUs)}
Gated Recurrent Units (GRUs) are a simplified version of LSTMs that also address the vanishing gradient problem. GRUs use only three gates:
\begin{itemize}
    \item \ul{Update Gate}: Controls how much of the previous information should be retained and how much of the new information should be added to the memory cell.
    \item \ul{Reset Gate}: Controls how much of the previous information should be ignored when computing the new candidate values.
    \item \ul{Candidate Gate}: Generates new candidate values that could be added to the memory cell.
\end{itemize}

\subsubsection{Minimal Gated Recurrent Units (M-GRUs)}
Minimal Gated Recurrent Units (M-GRUs) are an even simpler version of GRUs that use only two gates:
\begin{itemize}
    \item \ul{Update Gate}: Controls how much of the previous information should be retained and how much of the new information should be added to the memory cell.
    \item \ul{Candidate Gate}: Generates new candidate values that could be added to the memory cell.
\end{itemize}
